\chapter{Fattori chiave per l'utilizzo efficace dell'AI nello sviluppo software}

L'utilizzo di modelli di intelligenza artificiale all'interno del processo di sviluppo software non si limita alla semplice generazione automatica di codice. Durante lo sviluppo dell'applicazione realizzata in questa tesi è emerso come l'efficacia di tali strumenti dipenda in larga misura da come l'interazione con il modello viene strutturata. In particolare sono stati individuati tre fattori principali che influenzano in modo significativo la qualità dell'output generato: la formulazione dei prompt, l'organizzazione del contesto informativo fornito al modello e l'integrazione degli strumenti di AI all'interno del normale ambiente di sviluppo.

\section{Utilizzo di prompt efficaci}

Il primo fattore riguarda la formulazione dei prompt utilizzati per interagire con il modello. I modelli linguistici generativi non possiedono una comprensione diretta del progetto software su cui stanno operando, ma producono risposte basandosi esclusivamente sulle informazioni presenti nel prompt e nel contesto immediatamente disponibile. Di conseguenza, prompt generici tendono a produrre soluzioni vaghe, incomplete o non coerenti con l'architettura del sistema.
\\
Durante lo sviluppo dell'applicazione si è osservato che la qualità del codice generato migliora sensibilmente quando il prompt include informazioni esplicite riguardo ai vincoli tecnici del progetto. Tra gli elementi più rilevanti risultano in particolare l'indicazione del framework utilizzato, la descrizione della struttura del progetto e la specifica delle funzionalità che il codice deve implementare.
\\
Un ulteriore aspetto emerso riguarda la dimensione del problema assegnato al modello. I modelli di linguaggio tendono infatti a produrre risultati più affidabili quando vengono utilizzati per risolvere problemi circoscritti, come l'implementazione di una singola funzione o la modifica di una porzione limitata di codice. Al contrario, richieste troppo ampie o poco definite aumentano la probabilità di ottenere soluzioni incoerenti o non direttamente utilizzabili.
\\
Per questo motivo, durante lo sviluppo del progetto si è adottato un approccio incrementale, nel quale le richieste al modello venivano suddivise in compiti più piccoli e progressivi. Questa strategia consente di mantenere maggiore controllo sul codice generato e di verificare più facilmente la correttezza delle singole modifiche.

\section{Gestione del contesto e organizzazione dell'ambiente di sviluppo}

Un secondo fattore rilevante riguarda la gestione del contesto informativo disponibile al modello durante l'interazione. I modelli linguistici operano infatti su finestre di contesto limitate e non possiedono una conoscenza persistente dello stato del progetto. Quando il contesto fornito è insufficiente, il modello tende a generare soluzioni basate su assunzioni generiche o su pattern architetturali standard che possono non essere coerenti con la struttura effettiva dell'applicazione o generare ripetizioni inutili.
\\
Per ridurre questo problema, nel progetto sviluppato in questa tesi è stato adottato un approccio basato sull'utilizzo di file di documentazione testuale integrati direttamente all'interno del repository del progetto. In particolare sono stati introdotti diversi file Markdown e file di configurazione contenenti informazioni strutturate sull'architettura del sistema, sullo schema del database e sull'organizzazione dei moduli applicativi.
\\
Questi documenti vengono utilizzati come riferimento durante l'interazione con il modello e permettono di fornire un contesto più ampio rispetto al singolo file di codice su cui si sta lavorando. In questo modo l'intelligenza artificiale può generare soluzioni più coerenti con le scelte progettuali già adottate e ridurre la probabilità di introdurre modifiche incompatibili con l'architettura del sistema.
\\
Un ulteriore aspetto riguarda l'integrazione dell'AI all'interno dell'ambiente di sviluppo. L'utilizzo di strumenti integrati nell'editor consente di mantenere un accesso diretto ai file del progetto e di verificare immediatamente il comportamento del codice generato. In questo contesto il debugging e il testing del software assumono un ruolo fondamentale: il codice prodotto dal modello viene infatti validato tramite gli stessi strumenti utilizzati nello sviluppo tradizionale, permettendo di individuare rapidamente eventuali errori o incoerenze e di correggere iterativamente le soluzioni proposte.
